\SourceLeaf{369}%
auch nur bei der Kugelgestalt der Erde, und den Abstand einzelner Punkte auf der Erdfläche von einander zu berechnen, so wie überhaupt, dass der Gebrauch der sphärischen Trigonometrie nicht so ganz ausgeschlossen wäre. Auch bei der sphäroidischen Gestalt der Erde hätte wohl \textit{bestimmter} herausgehoben werden können, wie der Begriff der geographischen Breite anders modificirt werden müsse als auf der Kugel, und wie von dieser Breite die relative Lage gegen den Erdäquator, die Erdaxe und den Erdmittelpunkt abhängt. Doch diess sind Kleinigkeiten, die dem allgemeinen Werthe des Buchs keinen Abbruch thun, und auf die der Verfasser, wenn vielleicht eine neue Auflage erforderlich sein sollte, zu welcher ein für den Unterricht sehr empfehlenswerthes Buch wohl gelangen kann, leicht wird Rücksicht nehmen können.

\Needspace{7\baselineskip}
\pdfbookmark[0]{Anzeigen. Schwab und Metternich: Parallel-Linien. 1816.}{anzeigen-schwab-metternich-1816}
\begin{center}
\rule{.85\textwidth}{.3pt}\par
Göttingische gelehrte Anzeigen.\quad 1816 April 20.\par
\rule{.85\textwidth}{.3pt}
\end{center}
\hspace*{\parindent}\textit{Commentatio in primum elementorum} \textsc{Euclidis} \textit{librum, qua veritatem geometriae principiis ontologicis niti evincitur, omnesque propositiones, axiomatum geometricorum loco habitae, demonstrantur. Auctore} \textsc{J. C. Schwab}, \textit{Regi Württembergiae a consiliis aulicis secretioribus, academiae scientiarum Petropolitanae, Berolinensis et Harlemensis Sodali.} (65 Seiten in Octav.) Stuttgart 1814. Typis \textsc{J. F. Steinkopf}.

\textit{Vollständige Theorie der Parallel-Linien. Nebst einem Anhange, in welchem der erste Grundsatz zur Technik der geraden Linie angegeben wird. Herausgegeben von} \textsc{Matthias Metternich}, \textit{Doctor der Philosophie, Professor der Mathematik, Mitglied der gelehrten Gesellschaft nützlicher Wissenschaften zu Erfurt.} 44 Seiten in Octav. Mainz 1815. Auf Kosten des Verfassers in Commission bei \textsc{Florian Kupferberg}.

Es wird wenige Gegenstände im Gebiete der Mathematik geben, über welche so viel geschrieben wäre, wie über die Lücke im Anfange der Geometrie bei Begründung der Theorie der Parallel-Linien. Selten vergeht ein Jahr, wo nicht irgend ein neuer Versuch zum Vorschein käme, diese Lücke auszufüllen, ohne dass wir doch, wenn wir ehrlich und offen reden wollen, sagen könnten, dass wir im Wesentlichen irgend weiter gekommen wären, als \textsc{Euklides} vor 2000 Jah%
